\subsection{Обучение с учителем и без учителя, примеры}

Описанные выше задачи являются примерами задач обучения с учителем (supervised learning), так как правильные ответы
для каждого объекта обучающей выборки заранее известны.
Задачи обучения с учителем делятся на виды в зависимости от того, каким может быть множество $Y$ всех возможных
ответов (таргетов)

Есть и относительно небольшой класс задач, относящихся к обучению без учителя (unsupervised learning), -- это
задачи, для которых нам известны только данные, а ответы неизвестны или вообще не существуют.
Более того, часто поиск \("\)правильных\("\) ответов не является самоцелью.
Классическим примером обучения без учителя является кластеризация -- задача разделения объектов на группы,
обладающие некоторыми неизвестными нам, но, как мы в глубине души надеемся, интерпретируемыми свойствами.
Примером может служить кластеризация документов из электронной библиотеки по темам или кластеризация новостей с
целью выделения крупных сюжетов.
Бывают и другие виды (и даже парадигмы) машинного обучения

\subsection{Назвать другие категории}

Помимо основных, выделяют также
\begin{itemize}
    \item \textbf{Обучение с подкреплением (Reinforcement Learning)}: агент обучается, взаимодействуя со средой
    и получая награду за правильные действия (робот, играющий в шахматы, беспилотный автомобиль)
    \item \textbf{Частичное обучение (Semi-supervised Learning)}: используется небольшой набор размеченных и
    большой набор неразмеченных данных (классификация веб-страниц, распознавание речи)
    \item \textbf{Самообучение (Self-supervised Learning)}: метки генерируются автоматически из самих данных
    (предобучение больших языковых моделей, обучение представлений в компьютерном зрении)
    \item \textbf{Активное обучение (Active Learning)}: алгоритм может запрашивать разметку для наиболее
    информативных примеров у эксперта (разметка медицинских изображений, анализ юридических документов)
    \item \textbf{Трансферное обучение (Transfer Learning)}: использование знаний, полученных при решении одной
    задачи, для быстрого обучения на другой, но родственной задаче (дообучение ImageNet модели для конкретного
    класса изображений)
\end{itemize}

\subsection{Практические примеры (минимум 3 для каждого типа)}

Для обучения с учителем ниже.

\underline{Обучение без учителя}
\begin{enumerate}
    \item \textbf{Сегментация клиентской базы:} разделить всех клиентов интернет-магазина на группы по истории
    покупок и демографии для персонализации рекламы (кластеризация, например, методом k-means)
    \item \textbf{Обнаружение аномалий в работе сервера:} выявить нестандартные паттерны в логах сетевой
    активности, которые могут указывать на кибератаку (часто используется снижение размерности и анализ плотности)
    \item \textbf{Рекомендация статей пользователю:} на основе матрицы ``пользователь-просмотренная статья'' с
    помощью матричной факторизации (SVD) найти скрытые темы и рекомендовать не прочитанные, но релевантные статьи.
\end{enumerate}

\underline{Обучение с подкреплением}
\begin{enumerate}
    \item \textbf{Управление роботом-манипулятором:} агент (контроллер робота) учится совершать точные движения,
    чтобы взять предмет, получая положительную награду за успех и отрицательную за лишние действия или промах
    \item \textbf{Оптимизация энергопотребления в ЦОД:} система управления учится динамически регулировать
    охлаждение и нагрузку на серверы, чтобы минимизировать затраты на электроэнергию.
    \item \textbf{Торговый бот на бирже:} агент учится совершать операции покупки/продажи акций, максимизируя
    конечный капитал (награда) на исторических или симулированных данных.
\end{enumerate}

\subsection{Постановка практической задачи, основные сущности}

\begin{enumerate}

    \item $\mathbb{Y} = \mathbb{R}$ или $\mathbb{Y} = \mathbb{R}^M$ — регрессия.
    Примерами задач регрессии является предсказание продолжительности поездки на каршеринге, спрос на конкретный
    товар в конкретный день или погода в вашем городе на завтра (температура, влажность и давление -- это
    несколько вещественных чисел, которые формируют вектор нашего предсказания)

    \item $\mathbb{Y} = 0,1$ -- бинарная классификация.
    Например, мы можем предсказывать, кликнет ли пользователь по рекламному объявлению, вернёт ли клиент кредит
    в установленный срок, сдаст ли студент сессию, случится ли определённое заболевание у пациента, есть ли на
    картинке банан.

    \item $\mathbb{Y} = 1,\dots,K$ -- многоклассовая (multiclass) классификация.
    Например, определение предметной области для научной статьи (математика, биология, психология и так далее)

    \item $\mathbb{Y} = 0,1^K$ -- многоклассовая классификация с пересекающимися классами (multilabel classification).
    Например, задача автоматического проставления тегов для ресторанов (логично, что ресторан может одновременно
    иметь несколько тегов)

    \item $\mathbb{Y}$ — конечное упорядоченное множество → задача ранжирования.
    Основным примером является задача ранжирования поисковой выдачи, где для любого запроса нужно отсортировать все
    возможные документы по релевантности этому запросу.
    При этом оценка релевантности имеет смысл только в контексте сравнения двух документов между собой, её
    абсолютное значение информации не несёт

\end{enumerate}